\section{Introduction}
\IEEEPARstart{I}{ndustrial-agent} deployments must connect each task to its
production context, process signals, and equipment interfaces. Existing work
addresses end-to-end automation control, operator-reviewed control-sequence
generation, modular factory skills, and production scheduling
\cite{xia2025control,huang2026made,koehle2026llmdriven,waseem2026pretrained}.
Other studies target PLC code generation, context-aware control, resource
capability discovery, shop-floor coordination, digital-twin fault handling,
natural-language machine access, and robotic-line control
\cite{liu2026agents4plc,wu2025instructmpc,lim2025llmenabled,zhao2026shopfloor,gill2025llmdt,hofmann2025nlc}.
Across production lines, agent roles, signal maps, protocols, objectives, and
product/run context vary. A reusable framework must configure these
differences around shared task and operation services while keeping each
action tied to its active production context.

Consider two production scenes with different process signals and interfaces:
cast-film thickness control and annealing-zone temperature control. An agent
task in either scene needs production context, typed signal access, and a
route from a named operation to the configured device protocol. The framework
must also distinguish a command receipt from a subsequent process
observation. The shared task and operation services can therefore be
configured with scene-specific signal maps, protocols, and limits while
retaining a common execution path. In this paper, the two simulator scenes
demonstrate that configured composition; they do not measure the effort or
benefit of changing configurations.

This paper presents AgentWorkShop, a configurable industrial-agent framework
that connects collaborative tasks and registered harnesses to production
context, manufacturing execution system (MES) data, analysis and modeling
services, data-acquisition (DAQ) and write-control (DCW) nodes, and
protocol-specific drivers. Its access surfaces include command-line, web,
SDK, REST, MCP, and A2A interfaces. Task and chat routes enter the
collaboration runtime, whereas a direct REST tool-invocation route can call
the shared operation dispatcher without creating a task.
Scene blueprints define signal maps, objectives, and guards; channel profiles,
tool policies, and agent bindings scope operations. Agents and callers invoke
named operations through shared runtime checks; the framework does not choose
routes autonomously and retains proposals, command confirmation, and process
observations as distinct records. Fig.~\ref{fig:arch} presents the system
architecture, and Fig.~\ref{fig:runtime} shows the cross-module task flow.
Section~IV details the node-exploration and recipe lifecycles.

The paper makes three contributions:
\begin{enumerate}
\item It defines a cross-layer task-to-operation contract that resolves
  production context before an agent turn, rechecks each returned operation
  at a shared dispatcher, routes semantic process nodes to configured device
  drivers, and links command receipts separately from process observations.
\item It implements this contract through replaceable harness adapters,
  scene and production-context configuration, analysis and modeling
  services, semantic device access, and route-specific command and recipe
  lifecycles. The contribution is this composition and evidence boundary;
  it does not introduce a new agent-reasoning or process-control algorithm.
\item It evaluates two current-version simulator paths across cast-film and
  annealing configurations. The application-level run verifies protocol
  reads, bounded writes, readback, and data lineage, while both model
  candidates fail the configured gates. Six sequential tasks through one
  registered harness record bounded recipe trials, synthetic approval,
  verified readback, task completion, and run-tagged samples; all samples
  remain outside their target bands. The evidence supports these tested
  integration paths, but does not measure configuration-reuse benefit,
  successful process optimization, convergence, physical transfer, or
  cross-harness portability.
\end{enumerate}

Section~II reviews related industrial-agent and integration work.
Section~III presents the design rationale, architecture, and cross-module
execution flow; Section~IV explains module responsibilities and their
integrated operation. Section~V describes the experiments and results, and
Sections~VI and~VII discuss the evidence limits and conclude the paper.
